\begin{table}[H]
\centering
\caption{Productividad laboral departamental, 2014--2024pr}
\label{tab:dept_productividad_resumen_33}
\scriptsize
\begin{tabular}{lrrrr}
\toprule
Departamento & PIB/trab. 2024pr & Crec. & PIB/hora 2024pr & Crec. \\
\midrule
Quindío & 33,8 & 7,2\% & 15,1 & 8,2\% \\
Putumayo & 443,0 & 6,3\% & 202,1 & 6,6\% \\
Risaralda & 38,7 & 3,4\% & 16,4 & 4,1\% \\
Caquetá & 24,9 & 3,5\% & 10,5 & 3,9\% \\
Caldas & 36,6 & 2,4\% & 15,6 & 2,8\% \\
San Andrés y Providencia & 82,4 & 2,9\% & 33,3 & 2,8\% \\
Sucre & 22,0 & 2,0\% & 10,4 & 2,8\% \\
Bogotá D.C. & 67,4 & 2,1\% & 28,6 & 2,4\% \\
Boyacá & 47,5 & 1,5\% & 22,6 & 2,2\% \\
Tolima & 40,0 & 1,8\% & 16,9 & 1,8\% \\
Valle del Cauca & 54,5 & 1,4\% & 23,4 & 1,7\% \\
Santander & 62,0 & 1,4\% & 26,1 & 1,7\% \\
Chocó & 28,6 & 2,0\% & 11,9 & 0,9\% \\
Bolívar & 45,1 & 0,6\% & 19,8 & 0,8\% \\
Norte de Santander & 24,3 & 0,9\% & 10,0 & 0,8\% \\
Vichada & 92,2 & 1,3\% & 38,1 & 0,8\% \\
Nariño & 17,8 & -0,5\% & 9,1 & 0,5\% \\
Cauca & 26,6 & -0,2\% & 13,7 & 0,3\% \\
Córdoba & 22,0 & 0,2\% & 10,7 & 0,2\% \\
Antioquia & 47,7 & -0,1\% & 20,3 & 0,2\% \\
Amazonas & 50,2 & -0,1\% & 21,7 & 0,2\% \\
Magdalena & 23,9 & 0,0\% & 10,3 & 0,2\% \\
Guaviare & 39,1 & -1,0\% & 17,0 & -0,0\% \\
La Guajira & 24,8 & -0,7\% & 11,1 & -0,4\% \\
Arauca & 171,2 & -0,3\% & 66,8 & -0,8\% \\
Huila & 31,6 & -1,3\% & 14,2 & -0,9\% \\
Cesar & 32,6 & -1,6\% & 13,8 & -1,1\% \\
Guainía & 38,8 & -1,5\% & 17,0 & -1,2\% \\
Atlántico & 37,3 & -1,8\% & 16,0 & -1,3\% \\
Cundinamarca & 40,4 & -1,7\% & 17,4 & -1,5\% \\
Meta & 63,1 & -2,7\% & 25,9 & -2,1\% \\
Vaupés & 58,5 & -3,3\% & 26,2 & -2,7\% \\
Casanare & 165,2 & -3,3\% & 67,8 & -3,3\% \\
\bottomrule
\end{tabular}
\caption*{\footnotesize Nota: PIB por trabajador en millones de pesos constantes de 2015; PIB por hora en miles de pesos constantes de 2015. Fuente: cálculos propios con DANE y GEIH.}
\end{table}